\section{Resolución de ecuaciones no lineales}

\subsection{Motivación}
Nuestro objetivo es partir de una función:
\[
    f: [a,b]\subset\mathbb{R}\to\mathbb{R}
\]
Continua, y encontrar sus raíces, es decir, los puntos $\xi_i \in [a,b]$:
\[
    f(\xi_i) = 0
\]

\subsection{Raíces aisladas}

\begin{definición}[Raíz aislada]
    Se dice que $\xi$ es una raíz aislada de $f$ si existe un entorno $U$ de $\xi$ tal que $\xi$ es la única raíz de $f$ en $U$, es decir, si $\exists \epsilon > 0$ tal que:
    \[
        f(x) \neq 0 \quad \forall x \in ( \xi - \epsilon, \xi + \epsilon ) \setminus \{ \xi \}
    \]   
\end{definición}

Si $f$ es continua, sus raíces tienen que ser aisladas.